%============================================================
% 12. 实习 & 兼职打工
%============================================================
\section{实习 \& 兼职打工}\label{sec:work}

本章默认读者持有丹麦高等教育学生居留许可。\textbf{能否工作、能工作多少以及是否限特定雇主，最终以本人 SIRI 决定信和居留卡文字为准。}CPR、税卡或雇主愿意录用都不能替代工作许可。

%------------------------------------------------------------
\subsection{法规与时间限制}

\begin{table}[H]
  \small
  \rowcolors{2}{gray!10}{white}
  \centering
  \setlength{\tabcolsep}{6pt}
  %\caption{在读期间与毕业后求职签的工时限制}
  \label{tab:work-hours}
  \begin{tabularx}{\textwidth}{p{0.40\textwidth} p{0.22\textwidth} X}
    \toprule
    \textbf{时段} & \textbf{每月可工作时长} & \textbf{备注} \\
    \midrule
    9 月 – 翌年 5 月           & 不超过 90 小时 & 以居留卡上注明的工作权利为准 \\[2pt]
    6 / 7 / 8 月             & 可全职 & 上述工作许可持有人可在这三个月全职工作 \\[2pt]
    毕业后求职期 & 通常同上 & 获批的 6 个月或 3 年求职期通常附带有限工作许可 \\
    \bottomrule
  \end{tabularx}
\end{table}

\begin{grayleftbox}\textbf{建议: }
上述规则适用于就读国家认可高等教育项目并获相应工作许可的学生。自 2025 年 5 月 2 日起申请部分\textbf{非国家认可}高等教育项目居留者不再自动获得工作许可。务必查看居留卡上的具体权限；非法或超时工作可能导致警告、居留被撤销等后果。最新规则：\url{https://www.nyidanmark.dk/en-GB/Words-and-concepts/SIRI/Work-permits-for-students-in-higher-educational-programmes}
\end{grayleftbox}

\subsubsection{每月工时怎么记}

90 小时按\textbf{日历月}计算，不是按四周滚动，也不能把上月未用工时挪到下月。应自行保存合同、排班、工时系统截图和工资单，并把多个雇主的有薪工时合并核对。居留卡仍写旧的“每周 20 小时”者也应查看 SIRI 对 2024 年 7 月后规则的官方说明；有疑问时向 SIRI 取得书面答复。




%------------------------------------------------------------
\subsection{求职实习基本信息}

\begin{enumerate}[leftmargin=2em,label=\arabic*.]
  \item \textbf{在丹课程实习（实训）}\\
        Mandatory 或 optional internship 若属于课程组成部分，应先取得学校批准。若需要全职实习但最初学生许可未包含该权利，须先向 SIRI 申请相应全职工作许可，\textbf{获批前不能开始全职实习}。
  \item \textbf{境外实习 / 工作}\\
        需申请目标国实习签；长期离境可能 \textbf{注销丹麦居住身份}。
  \item \textbf{用工顺序}\\
        申请岗位前先确认自己的居留是否附带工作权利，并把允许的工时如实告知雇主。不要仅凭“学生身份”推定可以工作。
  \item \textbf{找工作前置}\\
        签约前雇主通常需要姓名、住址、CPR／税务信息和合法工作证明。丹麦银行账户并非签订劳动合同的统一法定前提；工资收款方式应与雇主确认，并及时办理税卡及本人名下的收款账户。
\end{enumerate}

\subsection{实习许可判断流程}

\begin{enumerate}[leftmargin=2.2em,itemsep=3pt]
  \item 向课程负责人确认实习是 mandatory 还是 optional、是否计 ECTS，以及学校能否书面批准岗位与任务。
  \item 检查本人原始 SIRI 决定信是否已经包含指定期间的全职实习许可；不要只看企业写的“internship”。
  \item 若未包含，按 SIRI 的 \emph{Work permit for study-related employment during studies} 页面准备学校批准和实习协议。协议应写明起止日期、工时和工作内容。
  \item 等待许可期间，只能在本人现有有限工作许可范围内工作；需要全职许可的实习不得提前全职开始。
  \item 实习在丹麦境外时，同时核对目的地签证/工作许可，以及长期离丹是否导致丹麦 CPR 或居留许可失效，见第~\ref{sec:residence-lifecycle}~节。
\end{enumerate}

SIRI 官方入口：\url{https://nyidanmark.dk/en-GB/You-want-to-apply/Study/Study-internship}

%------------------------------------------------------------
\subsection{如何寻找工作 / 实习}

\begin{table}[ht]
  \small
  \rowcolors{2}{gray!10}{white}
  \centering
  \setlength{\tabcolsep}{6pt}
  %\caption{常用求职渠道}
  \label{tab:job-channels}
  \begin{tabularx}{\textwidth}{p{0.5\textwidth} X}
    \toprule
    \textbf{渠道} & \textbf{用途} \\
    \midrule
    Jobindex / Jobnet / LinkedIn / Jobcenter & 主流招聘网站 \\
    院系 Career Day \& 校内 Job Bank & 适合寻找学生职位和实习 \\
    直接投递 & 企业未发布岗位亦可主动发送 Application \\
    \bottomrule
  \end{tabularx}
\end{table}

\begin{grayleftbox}\textbf{小贴士： }
学生居留、毕业求职居留和以工作为依据的居留是不同法律基础。转换身份或接受可能超出当前许可范围的工作前，应查阅 SIRI 的对应页面；不要假定新许可获批后所有旧权限会以同一日期自动衔接。前往其他国家实习或工作还须遵守目的地国家的签证和工作许可规定。
\end{grayleftbox}

%------------------------------------------------------------
\subsection{求职工具与技巧}

\subsubsection{CV 编写}
\begin{itemize}[leftmargin=1.5em]
  \item \textbf{开头}：联系方式 + 1–2 句与 JD 高度相关的个人简介。
  \item \textbf{照片}：丹麦雇主普遍接受简洁职业照。
  \item \textbf{精准投递}：避免海投，结合岗位关键字定制内容。
\end{itemize}

\subsubsection{面试注意点}
\begin{itemize}[leftmargin=1.5em]
  \item \textbf{肢体语言}：微笑、眼神交流。
  \item \textbf{细节}：手机静音、提前踩点。
  \item \textbf{Follow-up}：面试后可视公司文化发送简短感谢邮件（Dear + 名/姓）。
\end{itemize}

\begin{grayleftbox}\textbf{小贴士：}
公司岗位已满时，主动 \textit{speculative application} 也可能获得机会。多数人无法「一次面试即拿 offer」，可先接受过渡岗位，再继续冲击 dream job。
\end{grayleftbox}

\subsection{签约前和第一张工资单}

\small
\rowcolors{2}{gray!10}{white}
\begin{longtable}{p{0.27\linewidth}p{0.63\linewidth}}
  \toprule
  \textbf{项目} & \textbf{要核对的内容} \\
  \midrule
  \endfirsthead
  \toprule
  \textbf{项目} & \textbf{要核对的内容} \\
  \midrule
  \endhead
  \bottomrule
  \endfoot
  雇主与岗位 & 法定雇主名称、CVR、工作地点、岗位职责、开始日期、是否固定期限 \\
  工时与许可 & 每周/每月约定工时、排班、加班处理；与本人许可上限合并核对 \\
  工资与附加项 & 税前时薪/月薪、发薪频率、晚班/周末补贴、养老金、餐食或其他扣款 \\
  假期 & 带薪休假还是 holiday allowance；谁向 FerieKonto 报送、何时可申请 \\
  试用和终止 & 试用期、双方通知期限、病假报告方式、适用集体协议（如有） \\
  第一张工资单 & 工时、税前工资、AM-bidrag、预扣税、养老金、假期工资、实际到账账户 \\
\end{longtable}
\normalsize

雇佣至少一个月且平均每周超过 8 小时时，雇主依法应提供包含核心条件的雇佣合同；即使低于该门槛，也应争取书面确认。丹麦没有全国统一法定最低工资，工资和附加待遇常由集体协议或个人合同决定。官方说明：\url{https://lifeindenmark.borger.dk/working/work-rights/working-conditions/employment-contracts}

\subsection{工会、A-kasse 与求助边界}

工会（fagforening）和失业保险基金（A-kasse）不是同一机构，加入通常也不是自动的。工会可就合同、工资和工作条件提供成员服务；A-kasse 关系到符合条件时的失业保险。非欧盟学生和毕业生是否能领取某项福利，还会受到会员期、工作记录、居留许可及“不得领取特定公共福利”等规则影响。\textbf{不要仅因缴纳会费就假定可以领取补助}；加入前请分别向 A-kasse、SIRI 和必要时所在市镇取得书面确认。

%------------------------------------------------------------
%------------------------------------------------------------
\subsection{国际组织与专项资助机会}

联合国机构、欧盟机构、研究项目和中国国家留学基金委（CSC）的岗位、资助金额、对象和申请窗口每年会变。先从雇主官方招聘入口取得岗位说明和 offer，再到资助方当年度项目指南核对国籍、学籍、毕业年限、保险和派出手续；不要套用往年固定日期或金额。

\begin{itemize}[leftmargin=2em,itemsep=2pt]
  \item UN Careers：\url{https://careers.un.org}
  \item 国家留学基金管理委员会：\url{https://www.csc.edu.cn}
  \item 如岗位在丹麦境外，另查目的地工作/实习许可；如在丹麦全职实习，仍须完成本章的 SIRI 判断流程。
\end{itemize}


%------------------------------------------------------------
\subsection{其他打工渠道}

\begin{table}[ht]
  \scriptsize
  \rowcolors{2}{gray!10}{white}
  \centering
  \renewcommand{\arraystretch}{1.2}
  \setlength{\tabcolsep}{5pt}
  %\caption{部分打工类型与时薪}
  \label{tab:parttime_wage}
  \begin{tabularx}{\textwidth}{p{0.22\textwidth} X X}
    \toprule
    \textbf{类型} & \textbf{说明} & \textbf{申请时核对} \\
    \midrule

    \textbf{零售 / 商场} & 超市、服装店收银或理货 & 语言要求、晚班／周末班、集体协议及附加工资 \\
    \textbf{餐饮} & 后厨、配餐、服务等 & 合同工时、试工是否计薪、小费规则和食品卫生要求 \\
    \textbf{家政 / 清洁} & 家庭或商业场所清洁 & 雇佣主体、交通时间、保险和清洁用品责任 \\
    \textbf{送报纸} & 清晨派报；常见招聘关键词 \textit{avisbud} & 计薪方式、路线距离、车辆要求和恶劣天气安排 \\
    \textbf{平台配送} & 通过 APP 接单配送 & 雇佣／承包身份、税务、保险、车辆成本和最低工时承诺 \\
    \bottomrule
  \end{tabularx}
\end{table}

\begin{grayleftbox}\textbf{小贴士: }
\begin{itemize}
  \item 丹麦没有全国统一的法定最低工资；实际工资和附加待遇常由集体协议或个人合同决定。
  \item 接受工作前取得书面合同，核对税前时薪、工时、假期工资、养老金、解约条件和保险安排。
\end{itemize}
\end{grayleftbox}



%------------------------------------------------------------
\subsection{小结}

\begin{itemize}
  \item \textbf{量力而行}：课程压力大时谨慎接兼职；丹麦生活费虽高，但时间成本同样宝贵。
  \item \textbf{合法合规}：9 月至次年 5 月按日历月合计不超过 90 小时；6--8 月是否可全职仍以本人许可为准。
  \item \textbf{语言加分}：即便基础岗位，掌握基础丹麦语也能拓宽选择。
  \item \textbf{聚焦简历}：精准对应 JD，适当附照片，避免千篇一律。
\end{itemize}

\noindent\textit{本章官方信息核验日期：2026 年 9 月 2 日。}
